\chapter[Segitiga sebangun]{Segitiga sebangun}\label{chap:parallel}

Dua segitiga $A' B' C'$ dan $A B C$ disebut
\index{segitiga!segitiga sebangun}\index{segitiga sebangun}\emph{sebangun} (secara ringkas \index{30@$\sim$}$\triangle A' B' C'\z\sim\triangle A B C$) jika (1) sisi-sisi yang bersesuaian sebanding;
yaitu,
$$A' B'
=
k\cdot A B,
\quad
B' C'=k\cdot B C
\quad
\text{dan}
\quad
C' A'
=
k\cdot C A
\eqlbl{dist}
$$
untuk suatu koefisien $k>0$ (yang disebut \index{koefisien kesebangunan}\emph{koefisien kesebangunan}), dan (2) sudut-sudut yang bersesuaian sama, mungkin dengan perubahan tanda:
$$
\begin{aligned}
\measuredangle A' B' C'&=\pm\measuredangle A B C,
\\
\measuredangle B' C' A'&=\pm\measuredangle B C A,
\\
\measuredangle C' A' B'&=\pm\measuredangle CAB.
\end{aligned}
\eqlbl{angles}
$$

\parbf{Catatan.}
\begin{itemize}
\item Menurut \ref{thm:signs-of-triug}, pada ketiga persamaan di atas, semua tandanya dapat dianggap sama.

\item Persamaan-persamaan dalam \ref{dist} dapat ditulis ulang sebagai
\[\frac{A'B'}{AB}=\frac{B'C'}{BC}=\frac{C'A'}{CA}.\]

\item Jika $\triangle A' B' C'\sim\triangle A B C$ dengan koefisien $k=1$,
 maka $\triangle A' B' C'\z\cong\triangle A B C$.

\item Perhatikan bahwa ``$\sim$'' merupakan
\index{relasi ekuivalensi}\emph{relasi ekuivalensi}.
Artinya,
\begin{enumerate}[(i)]
\item $\triangle A B C\sim\triangle A B C$
untuk setiap $\triangle A B C$.
\item Jika $\triangle A' B' C'\sim\triangle A B C$, maka
$$\triangle A B C\sim\triangle A' B' C'.$$
\item Jika $\triangle A'' B'' C''\sim\triangle A' B' C'$ dan $\triangle A' B' C'\z\sim\triangle A B C$, maka
$$\triangle A'' B'' C''\sim\triangle A B C.$$
\end{enumerate}
\end{itemize}

\section{Kriteria kesebangunan}

Dengan menggunakan notasi baru ``$\sim$'', kita dapat menyatakan kembali Aksioma~\ref{def:birkhoff-axioms:4}:

\begin{thm}{Pernyataan ulang Aksioma~\ref{def:birkhoff-axioms:4}}
Jika untuk
$\triangle ABC$,
$\triangle AB'C'$,
dan $k>0$ berlaku
$B'\in [AB)$,
$C'\in [AC)$,
$AB'=k\cdot AB$ dan
$AC'=k\cdot AC$,
maka $\triangle ABC\sim\triangle AB'C'$.
\end{thm}

Dengan kata lain, Aksioma~\ref{def:birkhoff-axioms:4} memberikan
suatu syarat yang menjamin bahwa dua segitiga sebangun.
Mari kita rumuskan tiga {}\emph{kriteria kesebangunan} lainnya.

\begin{thm}{Kriteria kesebangunan}\label{prop:sim}
Dua segitiga
$\triangle ABC$ dan $\triangle A'B'C'$
sebangun jika salah satu kriteria berikut terpenuhi:

(SAS)\index{kriteria kesebangunan SAS}
Untuk suatu konstanta $k>0$ berlaku
\[A B=k\cdot A' B',
\quad
A C=k\cdot A' C',
\quad
\text{dan}
\quad
\measuredangle B A C=\pm\measuredangle B' A' C'.\]

(AA)\index{kriteria kesebangunan AA} Segitiga $A' B' C'$ tak degenerat
dan
$$\measuredangle A B C
=
\pm\measuredangle A' B' C',
\quad
\measuredangle B A C
=
\pm\measuredangle B' A' C'.$$

(SSS)\index{kriteria kesebangunan SSS} Untuk suatu konstanta $k>0$ berlaku
$$A B=k\cdot A' B',
\quad
A C=k\cdot A' C',
\quad
CB=k\cdot C'B'.$$

\end{thm}

Setiap kriteria ini dibuktikan dengan menerapkan Aksioma~\ref{def:birkhoff-axioms:4} bersama kriteria kekongruenan SAS, ASA, dan SSS secara berturut-turut
(lihat Aksioma~\ref{def:birkhoff-axioms:3} dan kriteria \ref{thm:ASA} serta \ref{thm:SSS}).


\parit{Bukti.}
Ambil $k=\tfrac{AB}{A'B'}$.
Pilih titik $B''\in [A'B')$ dan $C''\in [A'C')$,
sedemikian sehingga $A'B''=k\cdot A'B'$ dan $A'C''=k\cdot A'C'$.
Menurut Aksioma~\ref{def:birkhoff-axioms:4},
$\triangle A'B'C'\z\sim \triangle A'B''C''$.

Dengan menerapkan kriteria kekongruenan SAS, ASA, atau SSS, sesuai dengan kasusnya,
kita memperoleh $\triangle A'B''C''\cong \triangle ABC$.
Dengan demikian, pernyataan terbukti.
\qeds

\begin{thm}{Latihan}\label{ex:mid-triangle}
Misalkan $A'$, $B'$, dan $C'$ berturut-turut merupakan titik tengah sisi-sisi $[BC]$, $[CA]$, dan $[AB]$ dari $\triangle ABC$.
Tunjukkan bahwa $\triangle A'B'C'\sim\triangle ABC$ dan tentukan koefisien kesebangunannya.
\end{thm}

\begin{thm}{Latihan}\label{ex:k*triangle}
Misalkan $O$, $A$, $B$, $C$, $A'$, $B'$, dan $C'$ merupakan titik-titik yang berbeda.
Andaikan $A'\in [OA)$, $B'\in[OB)$, $C'\in [OC)$, $OA'=k\cdot OA$, $OB'=k\cdot OB$, dan $OC'=k\cdot OC$ untuk suatu $k>0$.
Tunjukkan bahwa $\triangle A'B'C'\sim\triangle ABC$.
\end{thm}

Suatu bijeksi $X\leftrightarrow X'$ dari bidang ke dirinya sendiri disebut \index{transformasi pelestari sudut}\emph{transformasi pelestari sudut} jika
\[\measuredangle ABC= \measuredangle A'B'C'\]
untuk setiap $\triangle ABC$ dan bayangannya $\triangle A'B'C'$.

(Istilah \index{transformasi}\emph{transformasi} digunakan untuk bijeksi bidang ke dirinya sendiri yang mempertahankan suatu struktur geometri tertentu.
Misalnya, {}\emph{transformasi isometrik} adalah \textit{transformasi yang mempertahankan jarak}.)


\begin{thm}{Latihan}\label{ex:angle-preserving-euclid}
Tunjukkan bahwa setiap transformasi pelestari sudut pada bidang mengalikan semua jarak dengan suatu konstanta tetap.
\end{thm}

\section{Teorema Pythagoras}\index{Teorema Pythagoras}

Ingat bahwa suatu segitiga disebut \index{segitiga!segitiga siku-siku}\index{segitiga siku-siku}\emph{siku-siku} jika salah satu sudutnya siku-siku.
Sisi yang berhadapan dengan sudut siku-siku disebut \index{hipotenusa}\emph{hipotenusa}.
Sisi-sisi yang mengapit sudut siku-siku disebut \index{sisi siku-siku}\emph{sisi-sisi siku-siku}.


\begin{thm}{Teorema}\label{thm:pyth}
Andaikan $\triangle ABC$ memiliki sudut siku-siku di~$C$.
Maka
$$AC^2+BC^2=AB^2.$$

\end{thm}

\parit{Bukti.}
Misalkan $D$ adalah titik kaki tegak lurus dari $C$ pada~$(AB)$.

\begin{wrapfigure}[4]{r}{40mm}
\vskip-2mm
\centering
\includegraphics{mppics/pic-66}
\end{wrapfigure}

Menurut Lema~\ref{lem:perp<oblique},
\begin{align*}
AD&<AC<AB,
\shortintertext{dan}
BD&<BC<AB.
\end{align*}
Oleh karena itu, $D$ terletak di antara $A$ dan $B$;
khususnya,
$$AD+BD=AB.\eqlbl{AD+BD=AB}$$

Menurut kriteria kesebangunan AA, kita memperoleh
$$\triangle ADC\sim\triangle ACB\sim \triangle CDB.$$
Khususnya,
$$
\frac{A D}{A C}=\frac{A C}{A B}
\quad
\text{dan}
\quad
\frac{B D}{B C}=\frac{B C}{B A}.
\eqlbl{BCD}$$

Mari kita tulis ulang kedua identitas dalam \ref{BCD}:
\begin{align*}
AC^2=AB\cdot AD
\quad
\text{dan}
\quad
BC^2=AB\cdot B D.
\end{align*}
Dengan menjumlahkan kedua identitas ini dan menerapkan \ref{AD+BD=AB}, kita memperoleh
$$AC^2 +BC^2=AB\cdot (AD+ B D)=AB^2.$$
\qedsf

Gagasan dalam bukti ini muncul dalam Elements \cite[X.33]{euclid},
tetapi bukti yang diberikan di sana \cite[I.47]{euclid} berbeda;
bukti tersebut menggunakan metode luas yang dibahas dalam Bab~\ref{chap:area}.


\begin{thm}{Latihan}\label{ex:pyth}
Andaikan $A$, $B$, $C$, dan $D$ seperti dalam bukti di atas.
Tunjukkan bahwa
$CD^2=AD\cdot BD$.

\end{thm}

Latihan berikut merupakan kebalikan Teorema Pythagoras.

\begin{thm}{Latihan}\label{ex:pyth-conv}
Misalkan $ABC$ merupakan segitiga sedemikian sehingga
$AC^2+BC^2=AB^2$.
Tunjukkan bahwa sudut di $C$ siku-siku.
\end{thm}

\section{Metode segitiga sebangun}

Bukti Teorema Pythagoras di atas menggunakan {}\emph{metode segitiga sebangun}.
Untuk menerapkan metode ini, kita perlu mencari pasangan-pasangan segitiga sebangun, kemudian menggunakan kesebandingan sisi-sisi yang bersesuaian dan/atau kesamaan sudut-sudut yang bersesuaian.
Pada awalnya, menemukan pasangan semacam itu mungkin tidak mudah.

{

\begin{wrapfigure}{r}{25mm}
\vskip-6mm
\centering
\includegraphics{mppics/pic-68}
\end{wrapfigure}


\begin{thm}{Latihan kelas}\label{ex:two-pairs-sim}
Misalkan $ABC$ merupakan segitiga tak degenerat, dan misalkan titik-titik $X$, $Y$, dan $Z$ terletak seperti pada gambar.
Andaikan $\measuredangle CAY\z\equiv\measuredangle XBC$.
Temukan empat pasang segitiga sebangun yang menggunakan keenam titik ini sebagai titik-titik sudut,
dan buktikan kesebangunannya.
\end{thm}

}

\begin{thm}{Latihan}\label{ex:ABC+D}
Andaikan $\triangle ABC$ tak degenerat.
Andaikan $D\in [AB]$ dan $\measuredangle CAB=\measuredangle BCD$.
Tunjukkan bahwa $CD=\tfrac{b\cdot a}c$, dengan $a=BC$, $b=CA$, dan $c=AB$.
\end{thm}

{

\begin{wrapfigure}{r}{25mm}
\vskip-4mm
\centering
\includegraphics{mppics/pic-69}
\end{wrapfigure}

\begin{thm}{Latihan}\label{ex:right-perp-bi}
Andaikan $\triangle ABC$ memiliki sudut siku-siku di~$C$.
Misalkan $M$ merupakan titik tengah $[AB]$.
Misalkan $P$ dan $Q$ berturut-turut menyatakan titik-titik potong garis sumbu $[AB]$ dengan sinar
$[BC)$ dan $[AC)$.
Tunjukkan bahwa
\[
4\cdot MP\cdot MQ = AB^2.
\]

\end{thm}


}

%\begin{thm}{Exercise}\label{ex:footpoints}
%Suppose that $\triange ABC$ and $\triange A'B'C'$ are nondegenerate,
%Assume $C'$ is a footpoint of $C$ on $(AB)$, $A'$ is a footpoint of $A$ on $(XC)$, $B'$ is a footpoint of $B$ on $(XC)$ for some point $X\in (AB)$
%Show that $\triangle ABC\sim \triange A'B'C'$.
%\end{thm}


\section{Ketaksamaan Ptolemy}

Suatu \index{segiempat}\emph{segiempat} didefinisikan sebagai susunan terurut empat titik berbeda di bidang.
Keempat titik ini disebut \index{titik sudut!segiempat}\emph{titik-titik sudut}.
Segiempat $ABCD$ juga akan dilambangkan dengan \index{25@$\square$}$\square ABCD$.

Untuk suatu segiempat $ABCD$,
keempat ruas $[AB]$, $[BC]$, $[CD]$, dan $[DA]$ disebut \index{sisi!segiempat}\emph{sisi-sisi $\square ABCD$};
dua ruas lainnya, $[AC]$ dan $[BD]$, disebut \index{diagonal!segiempat}\emph{diagonal-diagonal $\square ABCD$}.

\begin{thm}{Ketaksamaan Ptolemy}\label{ptolemy-inq}
Pada setiap segiempat, hasil kali kedua diagonal tidak dapat melebihi jumlah hasil kali pasangan-pasangan sisi yang berhadapan;
yaitu,
\[AC\cdot BD\le AB\cdot CD+ BC\cdot DA\]
untuk setiap $\square ABCD$.
\end{thm}

Kita akan menyajikan bukti klasik ketaksamaan ini dengan menggunakan metode segitiga sebangun beserta suatu konstruksi bantu.
Bukti ini diberikan sebagai ilustrasi --- bukti tersebut tidak akan digunakan lagi selanjutnya.

{

\begin{wrapfigure}{o}{33mm}
\vskip-3mm
\centering
\includegraphics{mppics/pic-70}
\end{wrapfigure}

\parit{Bukti.}
Perhatikan sinar $[AX)$ sedemikian sehingga $\measuredangle BAX=\measuredangle CAD$.
Dalam hal ini, $\measuredangle XAD\z=\measuredangle BAC$ karena jika $\measuredangle BAX$ ditambahkan pada ruas kiri dan $\measuredangle CAD$ pada ruas kanan, kedua ruas menjadi $\measuredangle BAD$.
Kita dapat mengambil
\[AX=\frac{AB}{AC}\cdot AD.\]

}

Dalam hal ini, berlaku
\begin{align*}\frac{AX}{AD}&=\frac{AB}{AC},
&
\frac{AX}{AB}&=\frac{AD}{AC}.\\
\shortintertext{Dengan demikian}
\triangle BAX&\sim \triangle CAD,
&
\triangle XAD&\sim\triangle BAC.\\
\shortintertext{Oleh karena itu}
\frac{BX}{CD}&=\frac{AB}{AC},
&
\frac{XD}{BC}&=\frac{AD}{AC},
\shortintertext{atau}
AC\cdot BX&=AB\cdot CD,
&
AC\cdot XD&=BC\cdot AD.
\end{align*}
Dengan menjumlahkan kedua persamaan ini, kita memperoleh
\[AC\cdot(BX+XD)=AB\cdot CD+BC\cdot AD.\]
Yang tersisa hanyalah menerapkan ketaksamaan segitiga, $BD\le BX+XD$.
\qeds

Dengan menggunakan bukti ini bersama \ref{prop:inscribed-quadrangle}, dapat ditunjukkan bahwa kesamaan hanya berlaku jika titik-titik sudut $A$, $B$, $C$, dan $D$ terletak pada suatu garis atau lingkaran dalam urutan siklik yang sama;
lihat juga \ref{ptolemy-id} untuk bukti lain bagi kasus kesamaan.
Latihan~\ref{ex:ptolemy} mengisyaratkan bukti lain untuk ketaksamaan Ptolemy dengan menggunakan koordinat kompleks.
